\chapter{Conclusions and Future Work}
\label{ch:conclusion}

This chapter brings the report to a close. Section~\ref{sec:conclusions} 
summarises the project's key contributions and findings in relation to the 
stated aims and objectives. Section~\ref{sec:future_work} outlines directions 
for future work that build on the results and limitations identified in 
Chapter~\ref{ch:evaluation}.

% ================================================================
\section{Conclusions}
\label{sec:conclusions}
% ================================================================

This project investigated how autonomous agents can learn cooperative and 
competitive football behaviour in a physically realistic multi-agent environment 
using deep reinforcement learning. The main aim was to design, implement, and 
evaluate a complete DRL pipeline progressing from foundational tabular methods 
to state-of-the-art policy optimisation, and to apply the latter to a MuJoCo 
Soccer 2v2 multi-agent task. Four measurable objectives guided the work.

To achieve these objectives, three experimental stages were implemented. 
First, Q-Learning with $\varepsilon$-greedy exploration was applied to the 
discrete FrozenLake-v1 environment, directly implementing the Bellman update 
rule. Second, a DQN with experience replay and a target network was implemented 
from scratch using PyTorch and applied to CartPole-v1, demonstrating neural 
function approximation. Third, PPO with parameter sharing, custom reward 
shaping, and the Stable-Baselines3 library was applied to the MuJoCo Soccer 2v2 
environment, trained on Google Colab with automatic cloud checkpointing.

The key findings of this project are as follows:

\begin{itemize}
    \item \textbf{O1 — Achieved.} The Q-Learning agent attained a 62.0\% win 
    rate on FrozenLake-v1, a 41-fold improvement over the 1.5\% random 
    baseline, confirming the convergence properties of the Bellman equation 
    on tabular problems.
    
    \item \textbf{O2 — Substantially achieved.} The DQN agent demonstrated 
    a clear and consistent upward learning trend on CartPole-v1, with the 
    rolling average score rising from 23.6 to 142.8 over 500 training 
    episodes, validating the experience replay and target network mechanisms.
    
    \item \textbf{O3 — Achieved.} The MuJoCo Soccer 2v2 environment was 
    successfully configured with Shimmy, PettingZoo, and SuperSuit, and 
    the custom three-component reward shaping wrapper was implemented and 
    verified to produce non-trivial learning signals from the first episode.
    
    \item \textbf{O4 — Achieved.} A shared PPO policy was trained on 
    Google Colab and evaluated over 10 episodes, demonstrating that the 
    four agents accumulate positive shaped rewards and exhibit purposeful 
    ball-chasing and goal-directed behaviour.
\end{itemize}

Overall, the central contribution of this project is a fully operational, 
cost-free multi-agent deep reinforcement learning pipeline for physically 
realistic football simulation, including reward shaping, parameter sharing, 
and fault-tolerant cloud training with automatic checkpointing. The results 
demonstrate that, even with a limited computational budget on free-tier GPU 
hardware, it is feasible to train DRL agents that learn meaningful cooperative 
behaviours in complex continuous environments.

% ================================================================
\section{Future Work}
\label{sec:future_work}
% ================================================================

The limitations identified in Chapter~\ref{ch:evaluation} suggest several 
concrete directions for future investigation:

\begin{itemize}
    \item \textbf{Extended training:} The most immediate improvement is 
    to train the PPO agent for 50–100 million timesteps using a paid Colab 
    subscription or a dedicated GPU server. Prior work in similar 
    environments~\citep{tassa2020dmcontrol} demonstrates that qualitatively 
    richer behaviours (structured passing, positional play) emerge only at 
    much larger training scales.
    
    \item \textbf{Reward coefficient tuning:} A systematic grid search 
    or Bayesian optimisation over the three reward shaping coefficients 
    ($\alpha_\text{dist}$, $\alpha_\text{vel}$, $\alpha_\text{goal}$) 
    would identify the combination that maximises learning speed and final 
    policy quality, replacing the current engineering-judgement values.
\end{itemize}

In the longer term, the following research directions could build on the 
outcomes of this project:

\begin{itemize}
    \item \textbf{Multi-Agent PPO (MAPPO) with centralised critic:} 
    Replacing the parameter-sharing architecture with a full CTDE framework 
    — where agents share a centralised value function during training but act 
    from local observations — would allow specialised role differentiation 
    (e.g., attacker vs.\ defender) and is expected to yield significantly 
    higher coordination quality~\citep{yu2022mappo}.
    
    \item \textbf{Self-play curriculum:} Introducing a self-play training 
    regime — where the opposing team is periodically replaced by a past 
    version of the learning team — would enable an automated difficulty 
    curriculum, progressively challenging the agents as their skill grows, 
    following the approach of~\cite{silver2016mastering}.
\end{itemize}

% ================================================================
\section{Summary}
\label{sec:con_summary}
% ================================================================

This chapter concluded the project by summarising the key findings in relation 
to the four stated objectives, all of which were achieved or substantially 
achieved within the available computational budget. The central contribution 
— a complete, cost-free, fault-tolerant multi-agent DRL pipeline for MuJoCo 
football — was established. Concrete short-term and long-term directions 
for future research were proposed, grounded in the limitations identified 
in Chapter~\ref{ch:evaluation}. Chapter~\ref{ch:reflection} provides a 
personal reflection on the project experience.